(1) 对平面定常绝热激波控制体，单位面积质量通量 $j=\rho_1u_1=\rho_2u_2$ 守恒。动量通量满足 $p_1+\rho_1u_1^2=p_2+\rho_2u_2^2$，无外功的总能量通量满足 $h_1+u_1^2/2=h_2+u_2^2/2$；理想气体比焓为 $h=\gamma p/[(\gamma-1)\rho]$。这些方程无需假设熵守恒。
(2) 由质量式得 $u_2=u_1/r$，将其代入动量和能量关系并消去 $p_2$，求得 $r=(\gamma+1)M_1^2/[(\gamma-1)M_1^2+2]$。再将压缩比代入动量方程，得到 $p_2/p_1=1+\frac{2\gamma}{\gamma+1}(M_1^2-1)$。对于 $M_1>1$，二者均大于一，流体在激波后减速、增压、密度升高。
(3) 理想气体状态方程给 $T_2/T_1=(p_2/p_1)/r$，故温度可由第二问的两比值显式求出。熵差为 $s_2-s_1=c_p\ln(T_2/T_1)-R_s\ln(p_2/p_1)$，其中 $R_s=c_p-c_v$ 是单位质量气体常数。对物理压缩支 $M_1>1$，该值为正；激波把有序动能不可逆地转化为内能，尽管能量通量守恒。
(4) $M_1\downarrow1$ 时 $r\to1$、$p_2/p_1\to1$、熵增趋零，激波退化为弱扰动。强激波 $M_1\to\infty$ 时 $r\to(\gamma+1)/(\gamma-1)$，密度压缩比有有限上界，而压力比按 $M_1^2$ 增长。若强行使用等熵关系 $p/\rho^\gamma=\mathrm{const}$，会错误消除熵产并给出不相容的压缩比。
